\begin{abstract}
Household food waste accounts for the largest share of global food
losses and continues to grow despite increased public awareness.
Existing smart-refrigerator and inventory-monitoring systems
typically optimise a single sensing modality, and few combine
product identification, environmental monitoring, and real-time
mobile interaction within one low-cost reference design. This paper
presents a cabinet-scale prototype that integrates four 50\,kg load
cells with an HX711 amplifier for shelf-weight estimation, a DHT11
sensor for temperature and humidity, an MQ135 sensor for indicative
gas observation, an AI Thinker ESP32-CAM for QR-based product
registration, and an LED illumination subsystem driven through a
MOSFET stage. Sensor output from an Arduino Uno is transported over
USB to a Python bridge that exposes the readings as a CORS-enabled
local HTTP endpoint, which is consumed by a cross-platform Flutter
application providing dashboard, camera, product, alert, and
configuration views. A tiled multi-QR routine recovers several
product codes from a single camera capture, and a pixel-level RGB
analysis yields an interpretable estimate of banana surface
spotting without dependence on machine-learning inference.
Demonstration sessions confirmed concurrent operation of all
subsystems with end-to-end latency below one second between the
microcontroller and the mobile dashboard. Limitations, particularly
with respect to lightweight-product detection and MQ135 baseline
drift, are documented alongside the qualitative results. The
complete firmware, bridge, and application source are publicly
archived for reproducibility.
\end{abstract}
